\documentclass{article} % For LaTeX2e
\usepackage{iclr2026_conference,times,fullpage}

% Optional math commands from https://github.com/goodfeli/dlbook_notation.
\input{math_commands.tex}
\input{new_commands.tex}

\usepackage{hyperref}
\usepackage{url}
\usepackage{xspace}
\usepackage[utf8]{inputenc}
\usepackage{amsmath}
\usepackage{amssymb}
\usepackage{listings}

\usepackage{xcolor}
\usepackage{array}
\usepackage{tikz}
\usetikzlibrary{patterns}
\usepackage{float}
\usepackage{booktabs}
\usepackage{algorithm}
\usepackage{algorithmicx}
\usepackage{algpseudocode}

\definecolor{commentteal}{RGB}{0, 128, 128}
\newcommand{\SG}[1]{\textcolor{commentteal}{[Saumya: #1]}}

\definecolor{commentpurple}{RGB}{128, 0, 128}
\newcommand{\JL}[1]{\textcolor{commentpurple}{[Julia Liu: #1]}}

\lstset{
    language=Python,
    basicstyle=\ttfamily\footnotesize,
    keywordstyle=\color{blue}\bfseries,
    commentstyle=\color{green!50!black}\itshape,
    stringstyle=\color{orange},
    showstringspaces=false,
    breaklines=true,
    frame=single,
    mathescape=true,
    morekeywords={function, each, yield, class, concat, gather, normalized, uniform, where, randn_like, scatter, clip_grad_norm_}
}

\usepackage[margin=1in]{geometry}

\title{Learning the Universe: Distribution-Matching for Cross-Fidelity Mapping in Astrophysical Simulations}
\author{Vedanti Kshirsagar, Julia Liu, Saumya Goyal, Matthew Ho, Barnabas Poczos}
\date{}

\begin{document}

\maketitle

\section{Introduction}

Extracting cosmological parameters from large-scale structure surveys means
comparing observations to simulations through simulation-based inference (SBI)
\cite{cranmer2020frontier}. SBI needs thousands of simulations to cover the
parameter space. High-fidelity simulations are too expensive for this, which
motivates fast emulators: explicit approximations such as the particle-mesh code
FastPM \cite{feng2016fastpm}, or generative neural emulators such as CHARM
\cite{pandey2025creating}, which produces high-resolution simulation outputs
directly from low-resolution inputs. Both introduce a distribution shift relative
to high-fidelity simulations. Left uncorrected, this shift biases downstream
inference.

This is a sim-to-real problem. \cite{bayer2025field} showed that different N-body
codes are out of distribution at the field level, and that Gaussian smoothing can
align them only by destroying small-scale information. Emulator shifts are usually
larger and more structured, because emulators approximate the full generative
process and not just the resolution. We need a correction that closes this gap
without throwing away cosmological information.

We pose this as a point cloud alignment problem. Simulation particles, halos, and
galaxies are naturally sets of points in phase space (position, velocity, mass,
structure). We seek a mapping $f_\psi$ from emulated catalogs to a distribution
compatible with high-fidelity simulations. We define compatibility by an
inference-consistency criterion: posteriors trained on transformed emulator output
should be unbiased relative to posteriors trained on high-fidelity data. This
task-driven notion of alignment, rather than pointwise reconstruction, is the core
contribution.

The framework is model-agnostic. The same pipeline applies to N-body particle
matching across gravity solvers, galaxy alignment across hydrodynamic codes, or any
pair of simulators that produce catalogs of the same physical objects. Here we
study the simplest case: halo-to-halo alignment between Quijote N-body simulations
and CHARM emulator outputs, run from shared initial conditions so that catalogs
correspond approximately object by object. We align object positions only.
Velocity, mass, structural features, and the fully unpaired setting are left to
future work.

\paragraph{Problem setup.}
We work with paired suites of halo catalogs spanning cosmological parameters
$\theta \in \mathbb{R}^m$ (for example $\Omega_m$ and $\sigma_8$). The high-fidelity
(HF) catalogs $\{\mathbf{X}_i\}$ come from the Quijote high-resolution latin
hypercube \cite{villaescusa2020quijote}. The low-fidelity (LF) catalogs
$\{\mathbf{X}'_i\}$ come from CHARM \cite{pandey2025creating}. The two suites have
different halo counts per catalog and per-halo features $d$ (position, velocity,
mass, structure), and outputs are permutation-invariant with respect to halo
ordering. Because both suites share initial conditions, halos correspond
approximately object by object. We voxelise each catalog onto a common real-space
grid, so each suite becomes a set of halo count fields. We seek a mapping $f_\psi$
from LF to HF count fields such that posteriors trained on $f_\psi(\mathbf{X}'_i)$
are unbiased relative to posteriors trained on $\{\mathbf{X}_i\}$ directly. CHARM
matches Quijote's spatial summary statistics to about $10\%$, but this residual
mismatch is a distribution shift that biases posteriors trained on
$\{\mathbf{X}'_i\}$ and evaluated on $\{\mathbf{X}_i\}$.

\section{Method}

We instantiate $f_\psi$ as a patch-wise, style-conditioned 3D generative
adversarial network. It maps an LF halo-count field to a super-resolved (SR) field
at the same spatial resolution.

\paragraph{Patching.}
Each $128^3$ simulation box is split into a $2\times2\times2$ grid of eight $64^3$
cores. The generator receives one $64^3$ LF core, with no extra neighbouring
context, and predicts a corrected core of the same size. The eight outputs tile the
full box exactly, so reconstruction is direct placement with no averaging or
overlap. Before entering the network, halo counts $n$ are transformed to model space
as
\[
    y = \log(1+n),
\]
and after inference the transform is inverted as
\[
    n = \exp(y)-1,
\]
with the result clamped to the physical output range.

\paragraph{Generator and discriminator.}
We use a constant-resolution StyleGAN2-style 3D convolutional network that learns
a same-resolution residual correction,
\[
    \mathrm{SR} = \mathrm{LF} + G_\psi(\mathrm{LF},s,z),
\]
where $s\in\mathbb{R}^5$ is the cosmology-conditioning vector
\[
    s=(\Omega_m,\Omega_b,h,n_s,\sigma_8),
\]
and $z$ denotes the injected noise fields. A $1\times1\times1$ style-modulated
convolution first projects the single-channel count field into a hidden feature
representation. Four constant-resolution convolutional blocks follow. Each block
has two style-modulated $3\times3\times3$ convolutions with periodic padding,
LeakyReLU activations, and learned-strength noise injection. Each block also has a
projection whose output contributes to the predicted residual. Style modulation
conditions the convolutional weights on $s$.

The discriminator $D_\phi$ is conditioned on the same vector $s$. It applies
style-modulated 3D convolutions with periodic padding and successive
average-pooling downsampling, followed by a final projection and spatial averaging
to produce one scalar logit per core.

\paragraph{Training objective.}
The generator is trained with a weighted combination of adversarial and voxel-wise
reconstruction losses,
\[
    \mathcal{L}_G
    =
    \lambda_{\mathrm{adv}}\mathcal{L}_{\mathrm{adv}}
    +
    \lambda_{\mathrm{rec}}\mathcal{L}_{\mathrm{rec}},
\]
where
\[
    \mathcal{L}_{\mathrm{adv}}
    =
    \mathbb{E}\!\left[
        \operatorname{softplus}\!\left(-D_\phi(\mathrm{SR},s)\right)
    \right],
    \qquad
    \mathcal{L}_{\mathrm{rec}}
    =
    \mathbb{E}\!\left[
        \|\mathrm{SR}-\mathrm{HF}\|_1
    \right].
\]
We use $\lambda_{\mathrm{adv}}=1$ and $\lambda_{\mathrm{rec}}=2$. The discriminator
is optimised with the matching non-saturating logistic GAN objective and a lazy
$R_1$ gradient penalty with $\lambda_{R_1}=10$, applied every 16 discriminator
updates. No power-spectrum loss is used in the current model; we return to this
choice in \S\ref{sec:pk}.

Both networks use Adam with learning rate $10^{-4}$ and
$(\beta_1,\beta_2)=(0,0.99)$, with the gradient norm clipped to $1$. Training runs
for 30 epochs over 1600 training simulations. One simulation, meaning its eight
$64^3$ cores, forms each optimisation batch.

\paragraph{Inference and stitching.}
At inference, the LF $128^3$ count field is transformed to model space and split
into eight cores. Each core is processed independently by $G_\psi$ with a fixed
seeded noise realisation. The cosmology $s$ is not known at deployment, so it is set
to a fiducial value $(0.30, 0.05, 0.70, 1.00, 0.80)$. The output cores are placed
directly into the full $128^3$ volume, and the stitched model-space field is
transformed back to physical counts.

\subsection{Baselines}

\textbf{Histogram regression.} A periodic U-Net (circular padding, four
$8{\to}16{\to}32{\to}64$-feature encoder and decoder stages) trained end to end with
a voxel-wise Poisson negative log-likelihood on the predicted log-rate. It has no
adversarial term, no patching, and no cosmology input, and it is applied to the
full $128^3$ box directly.

\textbf{LF + power modulation.} A non-learned correction. We fit a single
$k$-binned curve
$A(k)=\sqrt{\bar P_{\mathrm{HF,train}}(k)/\bar P_{\mathrm{LF,train}}(k)}$ on the
training split. For any LF field we then apply
$\delta(x)\to\mathrm{FFT}\to\times A(|k|)\to\mathrm{IFFT}\to
n(x)=\bar n\,(1+\delta'(x))$, clamped to be non-negative. Rescaling the Fourier
amplitude of a real field by a real, positive, isotropic $A(k)$ leaves its phase
unchanged, so this baseline's cross-coherence is identical to raw LF's by
construction,
\begin{equation}
    r_{\mathrm{powermod}}(k)
    =
    \frac{A(k)\,P_{\mathrm{LF,HF}}(k)}{\sqrt{A(k)^2 P_{\mathrm{LF}}(k)\,P_{\mathrm{HF}}(k)}}
    =
    r_{\mathrm{LF}}(k).
    \label{eq:powermod}
\end{equation}

\section{Results}

We test on the CMASS halo suite of $2000$ paired simulations at redshift $z=0.5$.
Each pair has a low-fidelity field (LF, from CHARM) and a high-fidelity field (HF,
from Quijote), run from the same initial conditions. Each field is a $128^3$ grid of
halo counts on a periodic box of side $L=1000\,\mathrm{Mpc}/h$. We split the suite
into $1600$ train, $200$ validation, and $200$ test simulations. Our model produces
a super-resolved field $\mathrm{SR}=f_\psi(\mathrm{LF})$. All power spectra are
computed on the count overdensity $\delta = n/\bar n - 1$.

\subsection{Metrics}

We measure three things: whether SR looks like HF (field statistics), whether SR
and HF fields come from the same distribution (a two-sample test), and whether SR
supports the same cosmology posterior as HF (the inference goal).

\paragraph{Power-spectrum statistics.} On $\delta$ we report the power spectrum
$P(k)$, the transfer function $T(k)=P_X(k)/P_\mathrm{HF}(k)$, and the
cross-coherence
$r(k)=P_{X,\mathrm{HF}}(k)/\sqrt{P_X(k)\,P_\mathrm{HF}(k)}$. $T(k)$ says whether a
field has the right amount of power at scale $k$, and is ideal at $1$. $r(k)$ says
whether the field has its structure in the right place, cell by cell, and is also
ideal at $1$. So $T(k)$ checks amplitude and $r(k)$ checks phase. Because every LF
box is paired with an HF box from the same initial conditions, we can also compute
the transfer per box and report its root-mean-square (RMS) departure from $1$. This
per-box error is a stricter test than the median transfer, and we use both.

\paragraph{PQMass.} To test whether two sets of fields share a distribution we use
PQMass \citep{lemos2024pqmass}. We pool the two sets and, many times over, pick
$n_r$ random reference points, assign every field to its nearest reference, and
count how many fields of each set land in each region. From these counts we form
the statistic
$\chi^2=\sum_i (\sqrt{N_2/N_1}\,k^{(1)}_i-\sqrt{N_1/N_2}\,k^{(2)}_i)^2/(k^{(1)}_i+k^{(2)}_i)$,
which follows a $\chi^2_{n_r-1}$ distribution if the two sets share a
distribution. If they differ, $\chi^2/\mathrm{dof}$ rises above $1$ and the
per-tessellation $p$-values move toward $0$. We describe each field by its
one-point count histogram, which is a property of the distribution and is not
dominated by the large-scale structure that all sources share through their common
initial conditions.

\paragraph{Cross-fidelity inference consistency (the main metric).} We train a
neural posterior estimator $q(\theta\mid s)$ on a summary $s$ of one source
$X\in\{\mathrm{HF},\mathrm{LF},\mathrm{SR}\}$ using the training split, then apply it
to HF test fields and measure how far its posterior sits from an HF-trained one. We
report the per-box Gaussian
$\mathrm{KL}\big(q_\mathrm{HF}(\theta\mid s_\mathrm{HF})\,\Vert\,q_X(\theta\mid s_\mathrm{HF})\big)$,
averaged over test boxes. A value of $0$ means source $X$ gives the same HF
posterior. Each estimator is a neural network trained with its own random seed, so
two estimators both trained on HF still disagree by a small amount. That HF-to-HF
value is the best score anything can reach, and we call it the floor. We evaluate
at two levels. At the summary level the input is $s=\log P(k)$, about $30$ numbers.
At the field level a 3D CNN reads the full $128^3$ field directly.

Two rules keep these numbers honest. First, every estimator is trained several
times with independent seeds and we report the mean and standard deviation across
seeds, not a single fit. Second, we check the HF-to-HF floor for every pair of
seeds. In our summary-level runs, one of five seeds produced a broken estimator
whose floor jumped from about $0.03$ to about $0.5$ against every other seed. That
seed is excluded from all summary-level numbers below. Roughly one in five
estimator fits fails this way, so no single cross-fidelity number should be quoted
to two decimals without this check.

\subsection{Field statistics: SR fixes the phase, not the amplitude}
\label{sec:pk}

Figure~\ref{fig:pkpanel} shows the per-patch power spectrum, transfer, and
cross-coherence for HF, LF, and SR. The transfer panel uses a linear axis centred on
$1$ and shows the 16th to 84th percentile band for both LF and SR, so a
ten-percent error is visible as a ten-percent offset.

\paragraph{Amplitude.} On the median, both LF and SR track the HF amplitude at all
scales. But the median hides the per-box error. Figure~\ref{fig:perbox} shows the
transfer per box, paired against the HF box from the same initial conditions, over
the $1600$ training boxes. The per-box RMS error of $P(k)$ against HF is $12.1\%$
for LF and $13.5\%$ for SR. SR does not reduce the error of the LF input, and it
beats LF on only $26\%$ of boxes. On the full $128^3$ box, the SR median transfer
sits at $0.90$ to $0.97$ for $k<0.1$, a systematic $5$ to $10\%$ deficit of
large-scale power, while LF sits slightly above $1$. Part of the per-box error is
high-$k$ shot noise that no model can remove, but SR removes none of the structured
part either. The cause is simple: the current model has no power-spectrum term in
its loss, so nothing pulls the per-box transfer toward $1$. The absolute RMS numbers
depend on the $k$-binning and the number of boxes, so we quote only the robust
statement: SR does not reduce the per-box power-spectrum error of LF.

\paragraph{Phase.} The picture is different for the cross-coherence. SR stays
correlated with HF at all scales, with $r(k)\gtrsim 0.5$ even near the Nyquist
frequency, while raw LF loses its correlation with HF at small scales and falls to
$r\approx 0$. SR puts small-scale structure back in the right place, which LF had
lost. This is the one field-level statistic where SR clearly improves on its input,
and it is the one that matters for the field-level inference result below.

\paragraph{Averaging samples does not help.} A natural idea is to average several
SR draws to form a posterior-mean field. We generated $10$ noise draws per patch
and averaged them (Fig.~\ref{fig:sampleavg}). The averaged field is identical to a
single draw: $T(k)$ and $r(k)$ change by less than $10^{-4}$ at every $k$. The
reason is that the generator is deterministic in practice. Its learned
noise-injection strengths are about $2\times10^{-4}$, compared with a
convolution-weight scale of about $0.02$, and the rounded output is bit-identical
across noise draws. We return to this in \S\ref{sec:pqmass}.

\begin{figure}[t]
\centering
\includegraphics[width=0.62\linewidth]{figures_cmass/pk_panel_patch_cmassAfix.png}
\caption{Per-patch power spectrum (top), transfer $T(k)=P_X/P_\mathrm{HF}$ on a
linear axis (middle), and cross-coherence $r(k)$ (bottom) for HF, LF, and SR, median
over $16$ test simulations with 16th to 84th percentile bands for both LF and SR.
The legend reports the RMS of $T(k)-1$. SR and LF have similar amplitude error, and
SR's scatter is no tighter than LF's. Only SR stays coherent with HF at small
scales. LF falls to $r\approx 0$.}
\label{fig:pkpanel}
\end{figure}

\begin{figure}[t]
\centering
\includegraphics[width=0.9\linewidth]{figures_cmass/pk_transfer_perbox_honest.png}
\caption{Per-box transfer against the paired HF box over $1600$ training
simulations. Left: median and 16th to 84th percentile band for LF and SR on a
linear axis. SR sits below $1$ for $k<0.1$. Right: histogram of the per-box RMS
power-spectrum error. SR ($13.5\%$) does not reduce the LF error ($12.1\%$).}
\label{fig:perbox}
\end{figure}

\begin{figure}[t]
\centering
\includegraphics[width=0.62\linewidth]{figures_cmass/pk_panel_patch_sampleavg_cmassAfix.png}
\caption{Same panel as Fig.~\ref{fig:pkpanel} with SR shown as a single noise draw
and as the mean over $10$ draws. The two SR curves lie exactly on top of each other
in all three panels: the generator has no sample diversity to average over.}
\label{fig:sampleavg}
\end{figure}

\subsection{Cross-fidelity: SR makes field-level inference transfer}

Table~\ref{tab:crossfid} and Fig.~\ref{fig:crossfid} give the main metric with
seed-to-seed error bars. We show three sources: raw LF, our deployed SR (the
cosmology-conditioned generator run with the fiducial $s$), and an unconditioned SR
variant trained with uninformative cosmology labels, which we use in the ablation
of \S\ref{sec:ablation}.

\paragraph{Field level.} A posterior trained on SR fields transfers to HF at or
near the floor, for both SR variants. A posterior trained on raw LF fields fails,
and fails unstably: its mean KL is more than ten times the floor and its standard
deviation exceeds its mean. Doubling the number of LF seeds did not shrink this
spread, so the instability is intrinsic. Some LF-trained networks transfer
acceptably and others break, with KL up to about $2$. The reason is that the raw LF
field is out of distribution for a field-level network, so the LF-trained posterior
gets the point estimate roughly right on HF but the posterior shape wrong. SR moves
the field into the HF distribution and makes the transfer reliable. This is the
main positive result: the correction unlocks field-level inference that raw LF does
not support.

\paragraph{Summary level.} On the power spectrum the story reverses. An LF-trained
summary posterior already sits close to the floor, and SR is worse. The ordering
LF $<$ SR is robust across seeds, although the SR spread is wide. Two facts explain
this. First, LF is already HF-like in the two-point amplitude that $P(k)$ measures
(\S\ref{sec:pk}), so there is little for the correction to fix at this level.
Second, SR does not just fail to help. It adds a per-box amplitude bias (the
large-scale deficit of Fig.~\ref{fig:perbox}) and, for the conditioned generator, a
cosmology-dependent distortion of $P(k)$ that misleads a posterior trained on it.
\S\ref{sec:ablation} separates these two effects.

\begin{table}[t]
\centering
\caption{Cross-fidelity KL to HF, lower is better, mean $\pm$ standard deviation
across independently trained estimators. Summary level: four healthy NDE seeds,
one broken seed excluded. Field level: all combinations of HF reference and
replicate. SR enables field-level transfer that raw LF cannot. At the summary level
LF is already close to the floor and SR is worse.}
\label{tab:crossfid}
\begin{tabular}{lcccc}
\toprule
level & LF (raw) & SR (deployed) & SR (unconditioned) & HF floor \\
\midrule
summary $P(k)$ & $\mathbf{0.049\pm0.004}$ & $0.160\pm0.036$ & $0.121\pm0.036$ & $0.030\pm0.005$ \\
field (3D CNN) & $0.47\pm0.78$ & $\mathbf{0.037\pm0.014}$ & $\mathbf{0.022\pm0.010}$ & $0.019$ \\
\bottomrule
\end{tabular}
\end{table}

\begin{figure}[t]
\centering
\includegraphics[width=0.9\linewidth]{figures_cmass/crossfid_honest_cmass.png}
\caption{Cross-fidelity KL to HF (log scale) at the summary and field levels, with
the HF-to-HF floor as the dashed line. Bars are the mean over seeds, error bars the
standard deviation, and dots the individual seeds at the summary level. SR reaches
the floor at the field level while raw LF fails with a spread larger than its mean.
At the summary level LF is near the floor and SR is worse.}
\label{fig:crossfid}
\end{figure}

\subsection{A distribution test, and why SR fails it}
\label{sec:pqmass}

PQMass on the one-point count histogram (Fig.~\ref{fig:pqmass}, $200$ validation
boxes) is a useful stress test. Raw LF cannot be told apart from HF
($\chi^2/\mathrm{dof}=0.83$, with $12\%$ of tessellations rejected at $p<0.05$), but
SR can be ($\chi^2/\mathrm{dof}=2.2$, $64\%$ rejected). So even though SR matches the
mean power spectrum and improves the cross-coherence, the distribution of its count
statistics differs from that of real halo counts. The gap is real but small: with
$60$ boxes SR sits inside the null bracket, and it takes about $200$ boxes to
separate it.

We ran two probes to find the cause. The first tests whether the problem is how
counts are produced (Fig.~\ref{fig:disc}). From one continuous generator output we
built three count fields: the raw continuous rate, the rounded field we deploy, and
a Poisson draw from the rate. Rounding changes nothing. The continuous and rounded
fields have the same one-point histogram (L1 distance to HF $0.0103$, versus
$0.0102$ for LF), the same PQMass score ($2.37$ versus $2.30$), and the same
transfer. Poisson sampling makes everything worse: PQMass rises to $30$, the
histogram distance grows tenfold, and the transfer climbs to $1.7$ at high $k$
because the draw adds white shot noise on top of a field that already contains the
clustered counts. The generator reproduces the count field directly, not a latent
rate, so a Poisson output head is the wrong fix and de-rounding is no fix at all.

The second probe measures the generator's sample variance
(Fig.~\ref{fig:samplevar}). Across $12$ noise draws per box the deployed SR field
does not change: the per-voxel standard deviation across draws is $0.00\%$ of the
mean, and the transfer band across draws is about $4000$ times thinner than the
band across boxes. The $L_1$ reconstruction term has driven the learned noise
strength to zero, so the generator outputs a single smoothed field per input.

Together these say the same thing. The mean one-point histogram of SR matches HF as
well as LF does, yet PQMass rejects it, so the mismatch is in the box-to-box spread
of the count statistics, not in their mean. A deterministic generator trained to a
mean-seeking loss is under-dispersed: it cannot reproduce the scatter that real
fields have from one realisation to the next. The lever is a distribution-matching
objective that restores stochastic scatter, not a different output head. This is
exactly the thesis of the paper, and the probes make the case for it directly.

\begin{figure}[t]
\centering
\includegraphics[width=0.9\linewidth]{figures_cmass/pqmass_Afix_fiducial.png}
\caption{PQMass two-sample test on the one-point count histogram, $200$ validation
boxes. Left: $\chi^2$ per tessellation against the null $\chi^2_{9}$ (dashed).
Right: $p$-values, which are uniform if the distributions match. LF matches HF;
SR is distinguishable.}
\label{fig:pqmass}
\end{figure}

\begin{figure}[t]
\centering
\includegraphics[width=\linewidth]{figures_cmass/discretization_Afix_fiducial.png}
\caption{Discretisation probe on $200$ validation boxes. From one continuous
generator output we form the continuous rate, the rounded field we deploy, and a
Poisson draw. Left: one-point count histogram against HF and LF. Middle: PQMass
$\chi^2/\mathrm{dof}$ against HF. Right: transfer to HF. Rounding changes nothing;
Poisson sampling adds shot noise and makes every metric worse.}
\label{fig:disc}
\end{figure}

\begin{figure}[t]
\centering
\includegraphics[width=0.9\linewidth]{figures_cmass/sample_variance_Afix_fiducial.png}
\caption{Generator sample variance, $60$ validation boxes and $12$ noise draws
each. Left: transfer with the band across draws (invisible) and the band across
boxes. Right: PQMass against HF for SR, LF, and the HF-to-HF floor at this sample
size. The generator is deterministic in practice, and at $60$ boxes SR is inside the
null bracket.}
\label{fig:samplevar}
\end{figure}

\subsection{Comparison to baselines}

Table~\ref{tab:baselines} and Fig.~\ref{fig:baseline_histreg} compare both
baselines with LF and HF on the test split. Histogram regression improves the
cross-coherence over raw LF at every scale we measure, the same qualitative gain
that $f_\psi$ makes, but it collapses small-scale power ($T(k)\approx0.03$ at
$k\approx0.29$, versus about $1$ for LF) and fails the PQMass test badly
($\chi^2/\mathrm{dof}=36.4$ versus $1.1$ for raw LF). Its output is a smooth Poisson
rate that is never resampled back to discrete counts, so its one-point histogram
lacks the shot noise that both LF and HF have. This is the same under-dispersion we
found in SR (\S\ref{sec:pqmass}), in a more severe form. Power modulation, by
contrast, is statistically indistinguishable from raw LF on PQMass
($\chi^2/\mathrm{dof}=1.10$), as expected for an amplitude-only transform that
barely touches the count distribution, and its coherence equals LF's exactly
(Eq.~\ref{eq:powermod}, Fig.~\ref{fig:baseline_powermod}).

The baselines sharpen the reading of Fig.~\ref{fig:pkpanel}. A transfer near $1$ is
not evidence of a good correction, because power modulation gets it by construction
and adds nothing. Raising the cross-coherence is not free, because histogram
regression gains coherence at a real cost in amplitude and one-point statistics.
$f_\psi$ is the only method of the three that improves coherence while keeping
$T(k)$ and the one-point histogram close to LF's. That is what the adversarial and
$L_1$ terms in $\mathcal{L}_G$ buy over a pointwise regression loss.

\begin{table}[t]
\centering
\caption{Non-generative baselines against raw LF on the $200$-simulation test
split, same field-statistics code as Fig.~\ref{fig:pkpanel}. $T(k)$ and $r(k)$ are
per patch at $k\approx0.03$ (large scale) and $k\approx0.29$ (small scale, near
Nyquist). PQMass uses the one-point count histogram with $n_r=10$, so
$\mathrm{dof}=9$. Note that the LF PQMass values here are on the test split and
differ slightly from the validation-split values in \S\ref{sec:pqmass}.}
\label{tab:baselines}
\begin{tabular}{lccccc}
\toprule
Method & $T(k{\approx}0.03)$ & $T(k{\approx}0.29)$ & $r(k{\approx}0.29)$ & PQMass $\chi^2/\mathrm{dof}$ & rejected ($p{<}0.05$) \\
\midrule
LF (raw)                    & $1.02$ & $1.02$ & $-0.11$ & $1.11$ & $19\%$ \\
Histogram regression        & $0.82$ & $0.03$ & $\mathbf{0.20}$ & $36.4$ & $100\%$ \\
LF + power modulation       & $1.01$ & $0.98$ & $-0.11$ & $1.10$ & $26\%$ \\
\bottomrule
\end{tabular}
\end{table}

\begin{figure}[t]
\centering
\includegraphics[width=0.62\linewidth]{figures_cmass/pk_panel_patch_histreg.png}
\caption{Histogram-regression baseline against HF and LF, per patch (median over
test patches, same convention as Fig.~\ref{fig:pkpanel}). The power spectrum tracks
HF out to $k\approx0.1$ and then collapses. The coherence exceeds raw LF's at every
scale.}
\label{fig:baseline_histreg}
\end{figure}

\begin{figure}[t]
\centering
\includegraphics[width=0.62\linewidth]{figures_cmass/pk_panel_patch_powermod.png}
\caption{LF + power-modulation baseline against HF and LF, per patch. The transfer
sits near $1$ by construction. The coherence curve lies exactly on top of raw LF's,
confirming Eq.~\ref{eq:powermod} numerically.}
\label{fig:baseline_powermod}
\end{figure}

\subsection{Ablation: conditioning the generator on cosmology}
\label{sec:ablation}

Since the goal is cosmological inference, it is natural to ask whether telling
$f_\psi$ the cosmology $\theta$ helps. We compare the conditioned generator with an
unconditioned variant trained on uninformative cosmology labels, and we run the
conditioned generator two ways: with the true $\theta$ of each box (an oracle, not
available at deployment) and with the fiducial value (the deployed setting).

Conditioning helps field fidelity when the true $\theta$ is supplied. The
validation power-spectrum RMS drops from $0.104$ (unconditioned) to $0.065$
(oracle). It does not help the deployment metric. At the field level all three
variants sit at or near the floor: $0.022$ unconditioned, $0.024$ oracle, $0.037$
fiducial. The fiducial value is slightly worse than the other two, which is what
one expects from feeding a fixed prior-mean cosmology to a generator that has
learned to use it, but the gap is within the seed-to-seed noise. At the summary
level conditioning hurts: $0.160\pm0.036$ for the deployed generator against
$0.121\pm0.036$ for the unconditioned one, and the oracle is no better than the
fiducial. The reason is that the power spectrum is the cosmology channel for a
summary posterior. A conditioned generator stamps its own imperfect mapping from
cosmology to $P(k)$ onto SR, and a posterior trained on SR then learns that mapping
instead of the true one. We conclude that a $\theta$-independent correction is the
better choice for cross-fidelity deployment, and that the value of the correction is
the small-scale phase structure it restores, not any cosmology it can be told.

\section{Next Steps}

\paragraph{Restore the power-spectrum loss.} The current model has no
power-spectrum term, and \S\ref{sec:pk} shows the consequence: SR carries a
coherent $5$ to $10\%$ large-scale power deficit and does not reduce the per-box
amplitude error of LF. A coherent deficit is exactly what a $P(k)$-matching term
fixes. We will re-add a per-patch power-spectrum loss and check that the per-box
transfer tightens toward $1$ without losing the coherence gain.

\paragraph{Restore stochasticity with a distribution-matching objective.}
\S\ref{sec:pqmass} shows the generator is deterministic and under-dispersed, and
that neither de-rounding nor a Poisson head fixes it. The $L_1$ term drives the
noise strength to zero because a mean prediction minimises it. We will replace or
reweight it with an objective that matches the distribution of fields rather than
the mean field, so that the box-to-box scatter of the count statistics matches HF.
This is the change the paper's thesis predicts, and PQMass gives a direct test of
whether it works.

\paragraph{Drop cosmology conditioning at deployment.} The ablation shows
conditioning helps only when the true $\theta$ is known, and hurts the summary
metric otherwise. The deployed model should be $\theta$-independent.

\paragraph{Redshift conditioning.} Right now we train a separate corrector for each
snapshot. Conditioning $f_\psi$ on redshift instead would let one model cover all
of them. Unlike cosmology, $z$ is always known when the model is used, so the
fiducial-value problem of \S\ref{sec:ablation} does not arise. We plan to train on
several snapshots, feeding $z$ through the same style pathway, and hold one redshift
out to test interpolation. We expect the LF to HF gap to be larger at low $z$, where
structure is more nonlinear, and will report cross-fidelity separately by redshift.

\paragraph{Patch stitching.} We measure $r(k)$ on individual patches, so it does not
capture seams created when patches are reassembled into the full $128^3$ box. The
field-level posterior reads the stitched field. Three changes should reduce seams:
generating overlapping patches and blending them with a smooth window, giving the
generator extra LF context around each patch, and wrapping patches at the box edge
to respect the periodic volume. Comparing $r(k)$ per patch against $r(k)$ on the
stitched field will measure how much stitching costs, and we will report it.

\bibliography{references}
\bibliographystyle{iclr2026_conference}

\end{document}
